\documentclass{article}
\usepackage[center]{titlesec}
\usepackage{graphicx} % Required for inserting images
\usepackage{tikz}
\usepackage{tikz-cd}
\usepackage{amsmath}
\usepackage{amssymb}
\usepackage{amsthm}
\usepackage[left=3cm, right=3cm, top=4cm, bottom=4cm]{geometry}
\usepackage{comment}
\usetikzlibrary{cd}
\usepackage{soul}
\usepackage{xcolor}
\usepackage{lmodern}
\usepackage[hidelinks]{hyperref}
\usepackage{cleveref}
\usepackage[enable]{darkmode}
\definecolor{marlight}{RGB}{196, 210, 0}
\sethlcolor{marlight}
\usepackage[style=numeric,backend=biber]{biblatex}
\addbibresource{references.bib}

\titleformat{\section}[block]{\normalsize\filcenter}{\thesection}{1em}{\MakeUppercase}


\numberwithin{equation}{section}

\theoremstyle{theorem}
\newtheorem{thm}{Theorem}[section]
\newtheorem{lem}[thm]{Lemma}

\theoremstyle{definition}
\newtheorem{defn}[thm]{Definition}

\theoremstyle{remark}
\newtheorem{rem}[thm]{Remark}

\theoremstyle{remark}
\newtheorem{exa}[thm]{Example}

\title{On Biadjoint Triangles}
\author{KB}
\date{\today}

\begin{document}

\pagenumbering{gobble}
\maketitle

\tableofcontents

\newpage

\pagenumbering{arabic}

\section{Monads}

\subsection{Monads and their Algebras}

It would be shameful not to begin our exploration without a good background in monads, which are described in detail in this section.
Monads, loosely speaking, answer the question of what it means for a category to possess some \textit{algebraic} structure.
This quality is made more explicit in the definition of \textit{monadicity}, which cements the idea that a category is \textit{monadic} over another category if it has some algebraic quality in relation to the latter category.
This intuition becomes all the more clear once examples of monads and monadicity are studied.
Much of this comes directly from Mac Lane, and will begin with a quick remark on composeability of functors and natural transformations.\\

For a category $\mathbb{X} $ and endofunctor $T:\mathbb{X} \to \mathbb{X}$, there are composite endofunctors $T^2 = T \circ T$ and $T^3 = T^2 \circ T$. 
If, additionally, there is a natural transformation $\mu: T^2\to T$, then, for each object $ x $ in $\mathbb{ X} $, the natural transformations $T\mu, \, \mu T: T^3\to T^2 $ have, respectively, components $(T\mu)_x= T\mu _x$, and $(\mu T)_x = \mu_{Tx} $ .

\begin{defn}\label{thm:monad-defn} \cite{MaclCT}
	A \textcolor{marlight}{monad} $T=( T,\eta, \mu ) $ on a category $\mathbb{X}$ is a triple in which $T:\mathbb{X}\to \mathbb{X}$ is an endofunctor on $ \mathbb{X} $ and $ \eta, \mu $ are two natural transformations
	\begin{equation}
		\eta:1_\mathbb{X}\longrightarrow T, \quad \mu: T^2\longrightarrow T
	\end{equation}
	such that the following diagram commutes:
	\begin{equation}\label{thm:monad-identities}
		% https://q.uiver.app/#q=WzAsNixbMCwwLCJUXjMiXSxbMSwwLCJUXjIiXSxbMiwwLCJUIl0sWzIsMSwiVF4yIl0sWzEsMSwiVCJdLFswLDEsIlReMiJdLFswLDEsIlxcbXUgVCJdLFsyLDEsIlxcZXRhIFQiLDJdLFszLDQsIlxcbXUiXSxbMSw0LCJcXG11IiwyXSxbNSw0LCJcXG11IiwyXSxbMCw1LCJUXFxtdSIsMl0sWzIsMywiVCAgXFxldGEiXSxbMiw0LCIiLDIseyJvZmZzZXQiOi0xLCJzdHlsZSI6eyJoZWFkIjp7Im5hbWUiOiJub25lIn19fV0sWzIsNCwiIiwyLHsic3R5bGUiOnsiaGVhZCI6eyJuYW1lIjoibm9uZSJ9fX1dXQ==
		\begin{tikzcd}
			{T^3} & {T^2} & T \\
			{T^2} & T & {T^2}
			\arrow["{\mu T}", from=1-1, to=1-2]
			\arrow["{T\mu}"', from=1-1, to=2-1]
			\arrow["\mu"', from=1-2, to=2-2]
			\arrow["{\eta T}"', from=1-3, to=1-2]
			\arrow[shift left, no head, from=1-3, to=2-2]
			\arrow[no head, from=1-3, to=2-2]
			\arrow["{T  \eta}", from=1-3, to=2-3]
			\arrow["\mu"', from=2-1, to=2-2]
			\arrow["\mu", from=2-3, to=2-2]
		\end{tikzcd}
	\end{equation}	
\end{defn}

Replace in \Cref{thm:monad-defn} the endofunctor $ T: \mathbb{X} \to \mathbb{X} $ with a set $ M $, $ T^2$ and $ T ^3 $ with the cartesian product $ M\times M $ and $  M\times M \times M $ respectively, and lastly the natural transformations $ \eta: 1_{ \mathbb{X}} \to T $ with $ \eta: 1 \to M $ and $ \mu: T^2 \to T $ with the binary operation $ \mu: M\times M \to M $.
Then the diagram in \Cref{thm:monad-identities} expresses the usual axioms of a monoid.
This is what motivates the naming of the natural transformations $ \eta $ as the \textcolor{marlight}{identity} and $ \mu $ as the \textcolor{marlight}{multiplication} of a monad $ (T, \eta, \mu) $.
This is also where we recover the famous ``a monad is just a monoid in \dots"\\

% TODO: Add something here

\begin{defn}
	Let $T=(T, \eta, \mu)$ be a monad on $ \mathbb{X}$.
	A \textcolor{marlight}{ $ T $-algebra} is a pair $(x,\xi)$ in which $x$ is an object in $ \mathbb{X}$ and $\xi: Tx \to x$ is a morphism in $ \mathbb{X} $ making the diagram
		\begin{equation}\label{def:t-algebra}
			% https://q.uiver.app/#q=WzAsNSxbMCwwLCJUXjIoWCkiXSxbMCwxLCJUKFgpIl0sWzEsMCwiVChYKSJdLFsyLDAsIlgiXSxbMSwxLCJYIl0sWzAsMiwiXFxtdV9YIl0sWzAsMSwiVChcXHhpKSIsMl0sWzMsMiwiXFxldGFfWCIsMl0sWzIsNCwiXFx4aSIsMl0sWzEsNCwiXFx4aSIsMl0sWzMsNCwiIiwwLHsic3R5bGUiOnsiaGVhZCI6eyJuYW1lIjoibm9uZSJ9fX1dLFszLDQsIiIsMCx7Im9mZnNldCI6LTEsInN0eWxlIjp7ImhlYWQiOnsibmFtZSI6Im5vbmUifX19XV0=
			\begin{tikzcd}
				{T^2x} & {Tx} & x \\
				{Tx} & x
				\arrow["{\mu_x}", from=1-1, to=1-2]
				\arrow["{T\xi}"', from=1-1, to=2-1]
				\arrow["\xi"', from=1-2, to=2-2]
				\arrow["{\eta_x}"', from=1-3, to=1-2]
				\arrow[no head, from=1-3, to=2-2]
				\arrow[shift left, no head, from=1-3, to=2-2]
				\arrow["\xi"', from=2-1, to=2-2]
			\end{tikzcd}
		\end{equation}
			
	commute. A \textcolor{marlight}{morphism of $T$-algebras} $h:(x,\xi) \to (x',\xi')$ is a morphism $h:x \to  x'$ in $ \mathbb{X}$ making the diagram below commute:
		\begin{equation}
			% https://q.uiver.app/#q=WzAsNCxbMCwwLCJUKFgpIl0sWzEsMCwiVChYJykiXSxbMCwxLCJYIl0sWzEsMSwiWCciXSxbMCwxLCJUKGgpIl0sWzAsMiwiXFx4aSIsMl0sWzEsMywiXFx4aSciXSxbMiwzLCJoIiwyXV0=
			\begin{tikzcd}
				{Tx} & {Tx'} \\
				x & {x'}
				\arrow["{Th}", from=1-1, to=1-2]
				\arrow["\xi"', from=1-1, to=2-1]
				\arrow["{\xi'}", from=1-2, to=2-2]
				\arrow["h"', from=2-1, to=2-2]
			\end{tikzcd}
		\end{equation}

\end{defn}

\begin{defn}
	For a monad $ T = (T, \eta, \mu) $ on $\mathbb{X} $, the \textcolor{marlight}{Eilenberg-Moore category} of $ T $, denoted by \textcolor{marlight}{$\mathbb{X}^T $}, is the category of $ T $-algebras with $ T $-algebra morphisms.
\end{defn}


\begin{thm}
	Every adjunction $ (F,U,\eta,\epsilon):\mathbb{X} \to \mathbb{A}$ determines a monad $ (UF, \eta, U \epsilon F) $ on $ \mathbb{X} $.
	
\end{thm}

\begin{proof}
	Firstly, $ T = UF $ is an endofucntor on $ \mathbb{X} $, $ \eta $ is a natural transformation $ \eta: 1_{ \mathbb{X}} \to UF =T $ and $ \mu = U \epsilon F: UFUF= T^2 \to T= UF$ is a natural transformation.
	For an object $ x $ in $ \mathbb{X} $, the commutativity of the diagram below must be checked:
		\begin{equation}\label{thm:proof-of-adj-monad}
			\begin{tikzcd}
				{UFUFUFx} & & {UFUFx} & & {UFx} \\
				& & & &\\
				{UFUFx} & & {UFx} & & {UFUFx}
				\arrow["{U\epsilon _{FUFx} }", from=1-1, to=1-3]
				\arrow["{UFU\epsilon _{Fx}}"', from=1-1, to=3-1]
				\arrow["U\epsilon _{Fx}"', from=1-3, to=3-3]
				\arrow["{\eta_{UFx}}"', from=1-5, to=1-3]
				\arrow[shift left, no head, from=1-5, to=3-3]
				\arrow[no head, from=1-5, to=3-3]
				\arrow["{UF \eta_x}", from=1-5, to=3-5]
				\arrow["U\epsilon _{Fx}"', from=3-1, to=3-3]
				\arrow["U\epsilon _{Fx}", from=3-5, to=3-3]
			\end{tikzcd}
		\end{equation}
	The left square commutes because of the naturality of $  \epsilon:FU \to 1_{ \mathbb{A}} $ applied to $ \epsilon_{Fx}:FUFx\to Fx $, which is preserved under $  U $.
	The right square commutes because of the triangle identities for the adjunction $ (F,U,\eta, \epsilon): \mathbb{ X} \to \mathbb{A} $.
	The identities are recalled:
		\begin{equation}\label{thm:triangle-identities-adjunction}
			% https://q.uiver.app/#q=WzAsNixbMCwwLCJGIl0sWzEsMCwiRlVGIl0sWzEsMSwiRiJdLFsyLDAsIlVGVSJdLFszLDAsIlUiXSxbMiwxLCJVIl0sWzAsMSwiRlxcZXRhIl0sWzEsMiwiXFxlcHNpbG9uIEYiXSxbMCwyLCIiLDIseyJvZmZzZXQiOi0xLCJzdHlsZSI6eyJoZWFkIjp7Im5hbWUiOiJub25lIn19fV0sWzAsMiwiIiwxLHsic3R5bGUiOnsiaGVhZCI6eyJuYW1lIjoibm9uZSJ9fX1dLFs0LDMsIlxcZXRhIFUiLDJdLFszLDUsIlUgXFxlcHNpbG9uIiwyXSxbNCw1LCIiLDAseyJvZmZzZXQiOjEsInN0eWxlIjp7ImhlYWQiOnsibmFtZSI6Im5vbmUifX19XSxbNCw1LCIiLDEseyJzdHlsZSI6eyJoZWFkIjp7Im5hbWUiOiJub25lIn19fV1d
			\begin{tikzcd}
				F & FUF & UFU & U \\
				& F & U
				\arrow["{F\eta}", from=1-1, to=1-2]
				\arrow[shift left, no head, from=1-1, to=2-2]
				\arrow[no head, from=1-1, to=2-2]
				\arrow["{\epsilon F}", from=1-2, to=2-2]
				\arrow["{U \epsilon}"', from=1-3, to=2-3]
				\arrow["{\eta U}"', from=1-4, to=1-3]
				\arrow[shift right, no head, from=1-4, to=2-3]
				\arrow[no head, from=1-4, to=2-3]
			\end{tikzcd}
		\end{equation}
	This is shown by first considering the bottom triangle in \Cref{thm:proof-of-adj-monad}, which is just the $ x $-component of the left triangle of \Cref{thm:triangle-identities-adjunction} under $ U $, and so it commutes.
	The top triangle in \Cref{thm:proof-of-adj-monad} is the $ Fx $-component of the right triangle in \Cref{thm:proof-of-adj-monad}.
	This shows that $  (UF, \eta, U \epsilon F) $ is a monad on $ \mathbb{X} $.
	This monad is called the \textcolor{marlight}{monad defined by the adjunction} $ (F,U, \eta, \epsilon): \mathbb{X}\to \mathbb{A} $.
\end{proof}

Then, naturally, the question arises of whether every monad induces an adjunction.
The answer is yes, and is discussed in the next theorem.

\begin{thm}
	Every monad $T=(T,\eta, \mu)$ on $ \mathbb{X}$ determines an adjunction $(F^T,U^T, \eta^T, \epsilon ^T): \mathbb{X} \to \mathbb{ X}^T$ from $ \mathbb{X}$ to the category of its $T$-algebras, $ \mathbb{X}^{T}$.
	The monad defined by this adjunction is the same monad $T=(T,\eta, \epsilon )$. 
	
\end{thm}

\begin{proof}
	Define
		\begin{enumerate}
			\item the right adjoint $U^{T}: \mathbb{X}^{T} \to \mathbb{X}$ to be the forgetful functor, $ U^{T}(x,\xi) := x$;
			\item the left adjoint $F^{T}: \mathbb{X} \to \mathbb{X}^{T}$ to be the free functor, $F^Tx:= (Tx, \mu_{x})$;
			\item the unit $\eta^{T}: 1_{ \mathbb{X}} \to U^{T}F^{T}$ to be the same $\eta$ of the monad;
			% TODO: very dodgy definition
			\item the counit $ \epsilon^{ T}: F^{T}U^{T} \to 1_{ \mathbb{X}^{T}} $ to be $ \epsilon^{ T}_{(Tx,\mu _x)} := \mu_x$. 
			
		\end{enumerate}		
	Now it is checked that this is indeed the desired adjunction.
	Firstly, note that $ (Tx,\mu_x) $ is a $ T $-algebra.
	That is, the diagram from \Cref{def:t-algebra} commutes:
		% https://q.uiver.app/#q=WzAsNSxbMCwwLCJUXjN4Il0sWzAsMSwiVF4yeCJdLFsxLDAsIlReMngiXSxbMiwwLCJUeCJdLFsxLDEsIlR4Il0sWzAsMiwiXFxtdV97VHh9Il0sWzAsMSwiVChcXHhpKSIsMl0sWzMsMiwiXFxldGFfe1R4fSIsMl0sWzIsNCwiXFxtdV94IiwyXSxbMSw0LCJcXG11X3giLDJdLFszLDQsIiIsMCx7InN0eWxlIjp7ImhlYWQiOnsibmFtZSI6Im5vbmUifX19XSxbMyw0LCIiLDAseyJvZmZzZXQiOi0xLCJzdHlsZSI6eyJoZWFkIjp7Im5hbWUiOiJub25lIn19fV1d
		\[\begin{tikzcd}
			{T^3x} & {T^2x} & Tx \\
			{T^2x} & Tx
			\arrow["{\mu_{Tx}}", from=1-1, to=1-2]
			\arrow["{T\mu_x}"', from=1-1, to=2-1]
			\arrow["{\mu_x}"', from=1-2, to=2-2]
			\arrow["{\eta_{Tx}}"', from=1-3, to=1-2]
			\arrow[no head, from=1-3, to=2-2]
			\arrow[shift left, no head, from=1-3, to=2-2]
			\arrow["{\mu_x}"', from=2-1, to=2-2]
		\end{tikzcd}\]
	But this is just the $ x $-component of \Cref{thm:monad-identities}, which commutes because $ T $ is a monad.\\
	
	Now to check naturality of $\eta^T$ and $\epsilon^T$.
	The naturality of $ \eta^T $ is just the naturality of $ \eta $.
	This is shown by taking $ x $ an object in $\mathbb{X} $, then the $ x $-component of $ \eta^T $ is the morphism in $ \mathbb{X} $, $ \eta^T_x: x \to U^TF^Tx = U^T(Tx, \mu_x)=Tx $, which is the same as $ \eta_x $.
	Now for $ \epsilon^T $: take $ (x, \xi) $ an object in $\mathbb{X}^T $. Then $ F^TU^T(x,\xi) = F^Tx = (Tx, \mu_x) $.
	Then the $ (x,\xi) $-component of $\epsilon^T$ is $ \epsilon^T_{(x,\xi)}:(Tx,\mu_x) \to (x,\xi) $.
	Taking a morphism of $ h:(x,\xi) \to (x',\xi') $ in $\mathbb{X}^T $ gives a naturality square:
		\begin{equation}\label{thm:naturality-of-epsilon-T}
			% https://q.uiver.app/#q=WzAsNixbMCwwLCIoeCxcXHhpKSJdLFswLDEsIih4JyxcXHhpJykiXSxbMSwwLCIoVHgsIFxcbXVfeCkiXSxbMSwxLCIoVHgnLFxcbXVfe3gnfSkiXSxbMiwwLCIoeCxcXHhpKSJdLFsyLDEsIih4JyxcXHhpJykiXSxbMCwxLCJoIiwyXSxbMiwzLCJGXlRVXlRoIiwyXSxbNCw1LCJoIl0sWzIsNCwiXFxtdV94Il0sWzMsNSwiXFxtdV94JyIsMl1d
			\begin{tikzcd}
				{(x,\xi)} & {(Tx, \mu_x)} & {(x,\xi)} \\
				{(x',\xi')} & {(Tx',\mu_{x'})} & {(x',\xi')}
				\arrow["h"', from=1-1, to=2-1]
				\arrow["{\mu_x}", from=1-2, to=1-3]
				\arrow["{F^TU^Th}"', from=1-2, to=2-2]
				\arrow["h", from=1-3, to=2-3]
				\arrow["{\mu_x'}"', from=2-2, to=2-3]
			\end{tikzcd}
		\end{equation}
		
	% TODO: must i elaborate on this?
	For some morphism $ f: x\to x' $ in $\mathbb{X} $, $ F^Tf: (Tx, \mu_x) \to (Tx', \mu_{x'}) $ must be a morphism of $ T $-algebras, which means $ F^Tf: Tx \to Tx' $ must be defined as $ F^Tf := Tf $.
	Then $ F^TU^Th = Th $. 
	So $ F^TU^Th $ in \Cref{thm:naturality-of-epsilon-T} becomes $ Th $, and the square thus commutes by naturality of $ \mu $, and becuase $ h $ is a morphism of $ T $-algebras.
	To check that there is an adjunction is to check the triangles in \Cref{thm:triangle-identities-adjunction} commute.
	For $ (x,\xi) $ an object in $\mathbb{X}^T $ and $ x $ and object in $ \mathbb{X} $, the left triangle is the $ (x,\xi) $-component of the left triangle of \Cref{thm:naturality-of-epsilon-T}, and the right is the $ x $-component of the right triangle of \Cref{thm:naturality-of-epsilon-T}:
		% https://q.uiver.app/#q=WzAsNixbMCwwLCJUeCJdLFsxLDAsIlReMngiXSxbMSwxLCJUeCJdLFsyLDAsIlR4Il0sWzMsMCwieCJdLFsyLDEsIngiXSxbMCwxLCJUXFxldGFfeCJdLFswLDIsIiIsMix7InN0eWxlIjp7ImhlYWQiOnsibmFtZSI6Im5vbmUifX19XSxbMCwyLCIiLDIseyJvZmZzZXQiOjEsInN0eWxlIjp7ImhlYWQiOnsibmFtZSI6Im5vbmUifX19XSxbMSwyLCJcXG11X3giXSxbNCwzLCJcXGV0YV94IiwyXSxbMyw1LCJcXHhpIiwyXSxbNCw1LCIiLDAseyJzdHlsZSI6eyJoZWFkIjp7Im5hbWUiOiJub25lIn19fV0sWzQsNSwiIiwxLHsib2Zmc2V0IjotMSwic3R5bGUiOnsiaGVhZCI6eyJuYW1lIjoibm9uZSJ9fX1dXQ==
		\[\begin{tikzcd}
			Tx & {T^2x} & Tx & x \\
			& Tx & x
			\arrow["{T\eta_x}", from=1-1, to=1-2]
			\arrow[no head, from=1-1, to=2-2]
			\arrow[shift right, no head, from=1-1, to=2-2]
			\arrow["{\mu_x}", from=1-2, to=2-2]
			\arrow["\xi"', from=1-3, to=2-3]
			\arrow["{\eta_x}"', from=1-4, to=1-3]
			\arrow[no head, from=1-4, to=2-3]
			\arrow[shift left, no head, from=1-4, to=2-3]
		\end{tikzcd}\]
	The left and right traingles commute respectively from the commutativity of the monad diagram \Cref{thm:monad-identities} and the $ T $-algebra diagram \Cref{def:t-algebra}.\\
	
	Lastly, notice that for an object $ x $ in $\mathbb{X} $, $ U^TF^Tx = Tx $, $ \eta^T = \eta $ and $ U^T\epsilon^TF^Tx=\mu_x $, and so the monad defined by this adjunction is $ T=(T,, \eta, \mu) $, as desired.		
	
\end{proof}

So we have that a monad determines an adjunction, which, in return determines the same monad again. The reverse is not always true: an adjunction determines a monad which does not always determine the same adjunction.
When an adjunction induces a monad that induces the same adjunction back, we call that adjunction \textit{monadic}. 
We get a sense of how monadic an adjunction is (or measure its \textit{monadicity} with the following construction.

\begin{thm}\label{thm:kiesli}
	Let $ (F,U,\eta,\epsilon):\mathbb{X}\to \mathbb{A}$ be an adjunction
	and $ T=(UF,\eta, U \epsilon F)$ be the monad defined by the adjunction. 
	Then there is a unique functor $K: \mathbb{A} \longrightarrow  \mathbb{X}^T $ with $U^TK= U$ and $KF=F^T$.
\end{thm}

% TODO: maybe actually prove this
% TODO: add stuff about desceeeeeent
% TODO: make a bibliography
\begin{proof}
	Define $K: \mathbb{A} \to \mathbb{X}^{T} $ by $Ka := (Ua, U \epsilon _a)$, and on morphisms $Kf:=Uf$.
\end{proof}

There is one last detail that will make our diagram complete: the functor $K$ as defined above has a left adjoint, the functor $L: \mathbb{X}^{T} \to \mathbb{A}$. 

\begin{thm}
	There is an adjunction $(L,K,\overline{\eta}, \overline{ \epsilon}): \mathbb{X}^{T} \to \mathbb{A}$.
\end{thm}

\begin{proof}
	We define
	\begin{enumerate}
		\item the left adjoint $L: \mathbb{X}^{T} \to \mathbb{A}$, where $L(x,\xi)$ is defined as the coequaliser of
		\[\begin{tikzcd}
			{FUFx} & {Fx} & {L(x,\xi)}
			\arrow["{\epsilon_{Fx}}", shift left, from=1-1, to=1-2]
			\arrow["{F\xi}"', shift right, from=1-1, to=1-2]
			\arrow["{\pi_{(x, \xi)}}", from=1-2, to=1-3]
		\end{tikzcd}\]
		\item the right adjoint $K: \mathbb{A} \to  \mathbb{X}^{T}$, defined by $Ka:=(Ua, U \epsilon_a)$ on objects and $Kf:=Uf$ on morphisms;
		\item the unit $\overline{\eta}: 1_{ \mathbb{X}^{T}} \to  KL$ defined by $\overline{\eta}_{(x,\xi)}=U\pi_{(x,\xi)}\eta_x$;
		\item the counit $\overline{ \epsilon}: LK \to 1_{ \mathbb{A}}$ is defined for each $ \overline{ \epsilon}_a$ as the unique morphism making the diagram
		\[\begin{tikzcd}
			{FUFUa} & {FUa} & {(LUa, U\epsilon_a)}\\ 
			& & {a}
			\arrow["{\epsilon_{FUa}}", shift left, from=1-1, to=1-2]
			\arrow["{FU\epsilon_a}"', shift right, from=1-1, to=1-2]
			\arrow["{\pi_{(Ua, U \epsilon_a)}}", from=1-2, to=1-3]
			\arrow["{ \epsilon_a}"', from=1-2, to=2-3]
			\arrow["{\overline{ \epsilon}_a}", from=1-3, to=2-3]
		\end{tikzcd}\]
		commute.
		
		
	\end{enumerate}
\end{proof} 

Now our picture is complete, and looks like below, with the $F$- and $U$-triangles commuting.
% https://q.uiver.app/#q=WzAsMyxbMiwwLCJcXG1hdGhiYntYfSJdLFswLDEsIlxcbWF0aGJie0F9Il0sWzIsMiwiIFxcbWF0aGJie1h9XlQiXSxbMSwwLCJVIiwwLHsib2Zmc2V0IjotMX1dLFswLDIsIkZeVCIsMix7Im9mZnNldCI6MX1dLFsyLDAsIlVeVCIsMix7Im9mZnNldCI6Mn1dLFswLDEsIkYiLDAseyJvZmZzZXQiOi0yfV0sWzEsMiwiSyIsMCx7Im9mZnNldCI6LTEsInN0eWxlIjp7ImJvZHkiOnsibmFtZSI6ImRhc2hlZCJ9fX1dLFsyLDEsIkwiLDAseyJvZmZzZXQiOi0yLCJzdHlsZSI6eyJib2R5Ijp7Im5hbWUiOiJkYXNoZWQifX19XSxbNiwzLCIiLDIseyJsZXZlbCI6MSwic3R5bGUiOnsibmFtZSI6ImFkanVuY3Rpb24ifX1dLFs0LDUsIiIsMix7ImxldmVsIjoxLCJzdHlsZSI6eyJuYW1lIjoiYWRqdW5jdGlvbiJ9fV0sWzgsNywiIiwyLHsibGV2ZWwiOjEsInN0eWxlIjp7Im5hbWUiOiJhZGp1bmN0aW9uIn19XV0=
\[\begin{tikzcd}
	&& {\mathbb{X}} \\
	{\mathbb{A}} \\
	&& { \mathbb{X}^T}
	\arrow[""{name=0, anchor=center, inner sep=0}, "F", shift left=2, from=1-3, to=2-1]
	\arrow[""{name=1, anchor=center, inner sep=0}, "{F^T}"', shift right, from=1-3, to=3-3]
	\arrow[""{name=2, anchor=center, inner sep=0}, "U", shift left, from=2-1, to=1-3]
	\arrow[""{name=3, anchor=center, inner sep=0}, "K", shift left, dashed, from=2-1, to=3-3]
	\arrow[""{name=4, anchor=center, inner sep=0}, "{U^T}"', shift right=2, from=3-3, to=1-3]
	\arrow[""{name=5, anchor=center, inner sep=0}, "L", shift left=2, dashed, from=3-3, to=2-1]
	\arrow["\dashv"{anchor=center, rotate=117}, draw=none, from=0, to=2]
	\arrow["\dashv"{anchor=center}, draw=none, from=1, to=4]
	\arrow["\dashv"{anchor=center, rotate=63}, draw=none, from=5, to=3]
\end{tikzcd}\]	
\\

\begin{defn} Let $ (F,U,\eta,\epsilon):\mathbb{X}\longrightarrow \mathbb{A}$ be an adjunction and $T= (T, \eta, \mu)$ be its associated monad. Then \begin{enumerate}
		\item the functor $ K: \mathbb{A} \longrightarrow \mathbb{X}^{T}$ as defined in the previous result is called the \textit{comparison functor};
		\item the functor $ U: \mathbb{A} \longrightarrow \mathbb{X}$ is said to be \textit{monadic} when the comparison functor is an equivalence.
	\end{enumerate}
	
\end{defn}



Monadic functors are very useful, because we can recover from them the adjunction that induces them. 
Moreover, monadicity means that, in some sense, a category is algebraic because it is equivalent to a category of algebras.\\


\subsection{Beck's Monadicity Theorem}

\begin{thm}\label{thm:beck}{\textit{Beck's Monadicity Theorem}}
	Let $ (F,U,\eta,\epsilon):\mathbb{X}\longrightarrow \mathbb{A}$ be an adjunction, and $T=(T,\eta, \mu)$ be its associated monad.
	The following conditions are equivalent:
	\begin{enumerate}
		\item the functor $ U: \mathbb{A} \longrightarrow \mathbb{X}$ is monadic;
		\item the functor $U$ preserves coequaliser diagrams of the form:
		% https://q.uiver.app/#q=WzAsMyxbMCwwLCJGVUYoWCkiXSxbMSwwLCJGKFgpIl0sWzIsMCwiTChYLFxceGkpIl0sWzAsMSwiXFxlcHNpbG9uX3tGKFgpfSIsMCx7Im9mZnNldCI6LTF9XSxbMCwxLCJGKFxceGkpIiwyLHsib2Zmc2V0IjoxfV0sWzEsMiwiXFxwaV97KFgsIFxceGkpfSJdXQ==
		\[\begin{tikzcd}
			{FUF(X)} & {F(X)} & {L(X,\xi)}
			\arrow["{\epsilon_{F(X)}}", shift left, from=1-1, to=1-2]
			\arrow["{F(\xi)}"', shift right, from=1-1, to=1-2]
			\arrow["{\pi_{(X, \xi)}}", from=1-2, to=1-3]
		\end{tikzcd}\]
		for every $T$-algebra $(X, \xi)$, and for every object $ A$ in $ \mathbb{A}$, the diagram
		\[\begin{tikzcd}
			{FUFU(A)} & {FU(A)} & {A}
			\arrow["{\epsilon_{FU(A)}}", shift left, from=1-1, to=1-2]
			\arrow["{FU(\epsilon_A)}"', shift right, from=1-1, to=1-2]
			\arrow["{ \epsilon_A}", from=1-2, to=1-3]
		\end{tikzcd}\]
		is a coequaliser;
		\item $U$ preserves coequaliser diagrams of the form:
		\[\begin{tikzcd}
			{FUF(X)} & {F(X)} & {L(X,\xi)}
			\arrow["{\epsilon_{F(X)}}", shift left, from=1-1, to=1-2]
			\arrow["{F(\xi)}"', shift right, from=1-1, to=1-2]
			\arrow["{\pi_{(X, \xi)}}", from=1-2, to=1-3]
		\end{tikzcd}\]
		for every $T$-algebra $(X,\xi)$, and $U$ reflects isos;
		\item every diagram 
		\begin{tikzcd}
			{A} & {B} 
			\arrow["{f}", shift left, from=1-1, to=1-2]
			\arrow["{g}"', shift right, from=1-1, to=1-2]
		\end{tikzcd}
		in $ \mathbb{A}$ - where
		\begin{tikzcd}
			{U(A)} & {U(B)} 
			\arrow["{U(f)}", shift left, from=1-1, to=1-2]
			\arrow["{U(g)}"', shift right, from=1-1, to=1-2]
		\end{tikzcd}
		has an absolute coequaliser -  	
		has a coequaliser preserved by $U$, and $U$ reflects isos.
		
	\end{enumerate}
	
\end{thm}

\begin{proof} Since it has been extensively covered that $(1)\iff (2) \iff(3)$, we will only prove $ (1)\implies (4)$.
	Let $ \mathcal{D}$ be a diagram in $ \mathbb{A}$ such that $ U( \mathcal{D}) =$ 
	\begin{tikzcd}
		{X} & {X'}
		\arrow["u", shift left, from=1-1, to=1-2]
		\arrow["v"', shift right, from=1-1, to=1-2]
	\end{tikzcd}
	with in $ \mathbb{X}$. Suppose $ U( \mathcal{D})$ has an absolute coequaliser $q$.
	
	% https://q.uiver.app/#q=WzAsNSxbMCwwLCJUKFgpIl0sWzEsMCwiVChYJykiXSxbMCwxLCJYIl0sWzEsMSwiWCciXSxbMiwxLCJRIl0sWzAsMiwiXFx4aSIsMl0sWzEsMywiXFx4aSciXSxbMCwxLCJUKHUpIiwwLHsib2Zmc2V0IjotMX1dLFswLDEsIlQodikiLDIseyJvZmZzZXQiOjF9XSxbMiwzLCJ1IiwwLHsib2Zmc2V0IjotMX1dLFsyLDMsInYiLDIseyJvZmZzZXQiOjF9XSxbMyw0LCJxIiwyXV0=
	\[\begin{tikzcd}
		{T(X)} & {T(X')} & \\
		X & {X'} & Q
		\arrow["{T(u)}", shift left, from=1-1, to=1-2]
		\arrow["{T(v)}"', shift right, from=1-1, to=1-2]
		\arrow["\xi"', from=1-1, to=2-1]
		\arrow["{\xi'}", from=1-2, to=2-2]
		\arrow["u", shift left, from=2-1, to=2-2]
		\arrow["v"', shift right, from=2-1, to=2-2]
		\arrow["q"', from=2-2, to=2-3]
	\end{tikzcd}\]	
	Since $ q$ is an absolute coequaliser of $ u$ and $ v$ in $ \mathbb{X}$, then $ T(q)$ is a coequaliser of $ T(u)$ and $ T(v)$ in $ \mathbb{A}$. 
	Since the $u$- and $ v$-diagrams commute, then $ q\xi'$ also coequalises $ T(u)$ and $ T(v)$. 
	Then there is a unique morphism $k: T(Q) \longrightarrow Q$ such that $ kT(q) = q \xi'$.
	Since this square commutes, then $ q$ is also a morphism in $ \mathbb{X}^{T}$.
	Now we want to show that $ q$ is the coequaliser of $ u$ and $ v$ in $ \mathbb{X}^{ T}$.
	Let $ q:(X',\xi') \longrightarrow (Q', k')$ be a morphism of $ T$-algebras such that $ q'u = q'v$ in $ \mathbb{X}^{T}$.
	Then $ q'u=q'v $ in $ \mathbb{X}$, and $ T(q')T(u)=T(q')T(v)$ in $ \mathbb{X}$, which, respectively, induce $l$ and $ T(l)$ in the diagram below, by universal property of the coequalisers.
	
	% https://q.uiver.app/#q=WzAsOCxbMCwwLCJUKFgpIl0sWzIsMCwiVChYJykiXSxbMCwyLCJYIl0sWzIsMiwiWCciXSxbMywyLCJRIl0sWzMsMCwiVChRKSJdLFs0LDEsIlQoUScpIl0sWzQsMywiUSciXSxbMCwyLCJcXHhpIiwyXSxbMSwzLCJcXHhpJyJdLFswLDEsIlQodSkiLDAseyJvZmZzZXQiOi0xfV0sWzAsMSwiVCh2KSIsMix7Im9mZnNldCI6MX1dLFsyLDMsInUiLDAseyJvZmZzZXQiOi0xfV0sWzIsMywidiIsMix7Im9mZnNldCI6MX1dLFszLDQsInEiXSxbMSw1LCJUKHEpIl0sWzEsNiwiVChxJykiLDIseyJsYWJlbF9wb3NpdGlvbiI6MzB9XSxbMyw3LCJxJyIsMl0sWzUsNCwiayIsMCx7InN0eWxlIjp7ImJvZHkiOnsibmFtZSI6ImRhc2hlZCJ9fX1dLFs2LDcsImsnIiwwLHsic3R5bGUiOnsiYm9keSI6eyJuYW1lIjoiZGFzaGVkIn19fV0sWzUsNiwiVChsKSJdLFs0LDcsImwiXV0=
	\[\begin{tikzcd}
		{T(X)} && {T(X')} & {T(Q)} & \\
		&&&& {T(Q')} \\
		X && {X'} & Q \\
		&&&& {Q'}
		\arrow["{T(u)}", shift left, from=1-1, to=1-3]
		\arrow["{T(v)}"', shift right, from=1-1, to=1-3]
		\arrow["\xi"', from=1-1, to=3-1]
		\arrow["{T(q)}", from=1-3, to=1-4]
		\arrow["{T(q')}"'{pos=0.3}, from=1-3, to=2-5]
		\arrow["{\xi'}", from=1-3, to=3-3]
		\arrow["{T(l)}", dashed, from=1-4, to=2-5]
		\arrow["k", dashed, from=1-4, to=3-4]
		\arrow["{k'}", from=2-5, to=4-5]
		\arrow["u", shift left, from=3-1, to=3-3]
		\arrow["v"', shift right, from=3-1, to=3-3]
		\arrow["q", from=3-3, to=3-4]
		\arrow["{q'}"', from=3-3, to=4-5]
		\arrow["l", dashed, from=3-4, to=4-5]
	\end{tikzcd}\]
	Now we just need to show that $l$ is also a morphism of $T$-algebras, which is to show that $k'T(l)=lk$. 
	Consider a segment of the diagram:
	% https://q.uiver.app/#q=WzAsNixbMCwwLCJUKFgnKSJdLFswLDEsIlgnIl0sWzEsMCwiVChRKSJdLFsxLDEsIlEiXSxbMiwwLCJUKFEnKSJdLFsyLDEsIlEnIl0sWzAsMiwiVChxKSJdLFsxLDMsInEiLDJdLFswLDEsIlxceGknIiwyXSxbMiwzLCJrIl0sWzIsNCwiVChsKSIsMCx7InN0eWxlIjp7ImJvZHkiOnsibmFtZSI6ImRhc2hlZCJ9fX1dLFszLDUsImwiLDIseyJzdHlsZSI6eyJib2R5Ijp7Im5hbWUiOiJkYXNoZWQifX19XSxbNCw1LCJrJyJdXQ==
	\[\begin{tikzcd}
		{T(X')} & {T(Q)} & {T(Q')} \\
		{X'} & Q & {Q'}
		\arrow["{T(q)}", from=1-1, to=1-2]
		\arrow["{\xi'}"', from=1-1, to=2-1]
		\arrow["{T(l)}", dashed, from=1-2, to=1-3]
		\arrow["k", from=1-2, to=2-2]
		\arrow["{k'}", from=1-3, to=2-3]
		\arrow["q"', from=2-1, to=2-2]
		\arrow["l"', dashed, from=2-2, to=2-3]
	\end{tikzcd}\]
	
	The outer square commutes because $ k' T(l)T(q)=k' T(q')=q'\xi'= lq \xi'$. 
	The left square also commutes by assumption, and T(q) is an epi since it is a coequaliser.
	Then by a pasting lemma, the right square commutes, too.
	Then $ l$ is indeed a $T$-algebra morphism, which shows that $ q$ is the coequaliser of $ u,v$ in $ \mathbb{X}^{T}$.
	From $ (1)$, $ \mathbb{X}^{T}\sim \mathbb{A}$, which means there is also a coequaliser of the diagram $ \mathcal{D}$ in $ \mathbb{A}$, which is trivially preserved by $ U$.
\end{proof}


Now we want a simpler way to check monadicity, which will become useful in the next example.

\begin{thm}\label{thm:crude}{ \textit{Crude Monadicity Theorem}}
	Let $ U: \mathbb{A} \longrightarrow \mathbb{X}$ be a functor. $ U$ is monadic if
	\begin{enumerate}
		\item $ U$ has a left adjoint;
		\item $ \mathbb{A}$ has, and $ U$ preserves, reflexive coequalisers;
		\item $ U$ reflects isomorphisms.
	\end{enumerate}
	
\end{thm}

\begin{proof}
	Let us only prove the second point. 
\end{proof}



\printbibliography

\end{document}
